\documentclass[11pt,a4paper]{article}
\usepackage[T1]{fontenc}
\usepackage{lmodern,amsmath,amssymb,booktabs,microtype}
\usepackage[margin=24mm]{geometry}
\usepackage[colorlinks=true,urlcolor=blue,linkcolor=blue]{hyperref}
\newcommand{\EE}{\mathbb E}
\newcommand{\core}{\mathcal E}
\DeclareMathOperator{\tr}{tr}
\DeclareMathOperator{\Rem}{Rem}
\title{An exact two-charge formula for the memory core\\Charged covariance transfer and its double-scaled filter}
\author{Alan Day \and Chengheng\thanks{AI research assistant (OpenAI).}}
\date{Local research note --- 6 November 2026}
\begin{document}\maketitle
\begin{abstract}
A charged-covariance transfer identity reduces the full quadratic
Gaussian contraction of the quartic memory core to two coherent
products in adjacent charge sectors. The resulting finite-$N$
coefficients are explicit, including the channels forbidden to
ordinary multiplication at half filling. After an exact cancellation,
the $H^2$ coefficient has a linear rank-deficit limit. Existing positive
product tails promote this part to an annealed $L^2$ approximation.
The weighted $A$ part and the output-deficit cross spectrum remain
the next steps toward the full memory observable.
\end{abstract}

\section{Conventions and the transfer identity}
Use standard CAR and the overlapping complex Gaussian $p$-body
ensemble, $\EE J_{IJ}J_{KL}=w\delta_{IL}\delta_{JK}$, $\hbar=1$.
The state is maximally mixed at fixed charge, with normalized trace
$\tr_m$ and $\|X\|_2^2=\EE\tr_m X^\dagger X$.
Set $h=H_m$, $k=H_{m-1}$, $c_i:\mathcal H_m\to\mathcal H_{m-1}$,
and $\Rem X=\sum_i c_i^\dagger Xc_i$. All finite formulas below assume
$m>p\ge1$; the double-scaled regime satisfies this condition.
Write $d=m-p$, $n=N-m+1$, and
\begin{align}
 C&=w\binom mp\binom{N-m+p}p,&
 C_-&=w\binom{m-1}p\binom{N-m+p+1}p,\\
 j_c&=w\binom{m-1}p\binom{N-m+p}p,&
 \kappa_r&=w\binom{m-r}p\binom{N-m+p-r}p,\\
 e_r&=\frac{w}{p}\binom{m-r-1}{p-1}
                  \binom{N-m+p-r}{p-1}.
\end{align}
Out-of-range binomials are zero. In the interior,
$e_r=p\kappa_r/[(m-r)(n-r)]$.

For an independent copy $(u,v)$ sharing the same new couplings in
the two sectors, define $\Psi(X)=\EE uXv$. For arbitrary rectangular
matrices $Y_i:\mathcal H_{m-1}\to\mathcal H_m$, the key identity is
\begin{equation}
 \boxed{P_r\sum_i\Psi(Y_i)c_i
 =\kappa_r P_r\sum_iY_ic_i
       -e_rP_r\Rem\!\left(\sum_i c_iY_i\right).}
 \label{eq:transfer}
\end{equation}
No rank restriction on $Y_i$ is required. The output rank is
$0\le r\le\min(m,N-m)$.

\newpage
\section{Why only two charged components are needed}
Let $T_r^\pm$ be the same primitive traceless rank tensor restricted
to charges $m$ and $m-1$. Introduce vector-valued operators
\begin{equation}
 V_i=T_r^+c_i^\dagger,\qquad U_i=c_i^\dagger T_r^-.
\end{equation}
Their difference is a commutator with one creation operator; it is
the traceless charged tensor of type $(r,r-1)$. The remaining
component has type $(r+1,r)$. Normal ordering with a single extra
creation operator permits precisely these two types. Their known
embedded-covariance eigenvalues are
\begin{align}
 L_r&=w\binom{m-r}p\binom{N-m+p-r+1}p,\\
 U_r^{\rm eig}&=w\binom{m-r-1}p\binom{N-m+p-r}p.
\end{align}
Put $D_r=N+1-2r$. Then
\begin{equation}
 L_r-U_r^{\rm eig}=D_r e_r,\qquad
 U_r^{\rm eig}+(n-r)e_r=\kappa_r.
\end{equation}
Use unnormalized HS norms in this paragraph, and set
$A_+=\|T_r^+\|^2$, $A_-=\|T_r^-\|^2$. Removal and addition give
\begin{equation}
 \frac{A_-}{A_+}=\frac{m-r}{n-r},\quad
 \|V-U\|^2=\frac{rD_r}{n-r}A_+,\quad
 \langle V,V-U\rangle=rA_+.
\end{equation}
Thus the low component of $V$ is $(n-r)(V-U)/D_r$, and
\begin{equation}
 \Psi(V_i)=\kappa_r V_i-(n-r)e_rU_i.
\end{equation}
At $r=0$ the difference is zero; endpoints follow by the surviving
component or the binomial formula.

Define $\mathsf S(Y)=\sum_iY_ic_i$ and
$\mathsf T(Y)=\sum_i c_iY_i$. Their HS adjoints are
$\mathsf S^*(X)_i=Xc_i^\dagger$ and
$\mathsf T^*(Z)_i=c_i^\dagger Z$. Also
$\sum_i c_iT_r^+c_i^\dagger=(n-r)T_r^-$.
Taking the adjoint of the preceding eigenvector identity, using
self-adjointness of $\Psi$, proves Eq.~\eqref{eq:transfer} on every
$Y$, including all product multiplicities.

For the embedded Hamiltonian ranks $h_a,k_a$, the useful consequences
are
\begin{align}
 \Psi(h_ac_i^\dagger)&=\kappa_a h_ac_i^\dagger
                   -\frac{p\kappa_a}{d}c_i^\dagger k_a,\\
 \Psi(c_i^\dagger k_a)&=-de_a h_ac_i^\dagger
                   +[\kappa_a+(m-n)e_a]c_i^\dagger k_a,\\
 \kappa_a(m-1)&=\kappa_a+(m-n-p)e_a.
\end{align}
These preserve the same couplings across the two charges.

\newpage
\section{The complete finite-$N$ quadratic core}
Recall
\begin{equation}
 \core=\sum_i(\operatorname{ad}_H^2c_i)^\dagger
                \operatorname{ad}_H^2c_i+\tfrac14[H,[H,A]],
 \qquad Q_{\core}=\tfrac12\Delta\core.
\end{equation}
The old exact contraction is
$Q_{\core}=D_4-2D_{13}+5D_{22}/4+7D_{121}/2+(4p-3m)D_{40}$,
with $D$ the contractions of the five raw words.

For $a<b$ use $X\odot Y=\{X,Y\}$; for $a=b$ use $X\odot X=X^2$.
Set
\begin{equation}
 Z^+_{ab;r}=P_r(h_a\odot h_b),\quad
 Z^-_{ab;r}=P_r\Rem(k_a\odot k_b),\quad
 x_a=a(N+1-a).
\end{equation}
Then, with all sums over allowed $a\le b,r$,
\begin{equation}
 \boxed{Q_{\core}=\sum_{a\le b,r}
          \left(g^+_{ab;r}Z^+_{ab;r}+g^-_{ab;r}Z^-_{ab;r}\right).}
 \label{eq:core}
\end{equation}
The coefficients are
\begin{align}
 g^+={}&(8p-2m)C_-+\frac{19}{2}d j_c+(12p-9m)C\notag\\
       &+m(\kappa_a+\kappa_b)+\frac{3p-m}{2}\kappa_r
                  -\frac{11}{4}dn(e_a+e_b),\label{eq:gp}\\
 g^-={}&3C_-+\frac52 C-4j_c
        -\frac{4m+p}{4d}(\kappa_a+\kappa_b)+\frac12\kappa_r\notag\\
       &-(m-n+p)(e_a+e_b)
        +\left[m+3n-p-\frac{5(2mn-x_a-x_b)}{4d}\right]e_r.
 \label{eq:gm}
\end{align}
No division by $Z^+$ occurs. The previously found $(2,2;1)$ channel
at $N=6,m=3$ is consequently retained.

To see the reduction, put $v_a=\kappa_a-ne_a$ and
$u_a=d(j_c+v_a)$. The old linear $U$ is $\sum_a u_ah_a$.
For its apparently new quadratic term, Eq.~\eqref{eq:transfer} gives
\begin{equation}
 P_r\sum_i\Psi(h_ac_i^\dagger k_b)c_i
 =d\kappa_rP_rh_ah_b
  -\frac{e_r(m-a)(n-a)}d P_r\Rem(k_ak_b).
\end{equation}
All other terms use the three consequences on the preceding page.
Adding their Hermitian conjugates produces the symmetric products
and yields Eqs.~\eqref{eq:gp}--\eqref{eq:gm}.

Gaussian centering acts after this identity:
$W_2\core=\sum(g^+W_2Z^++g^-W_2Z^-)$, since
$W_2\core=Q_{\core}-2\EE\core$.
The boundary $m=p$ is separate: $k=0$, $\core=p h^4$, and
$Q_{\core}=pD_{40}$.

\newpage
\section{The finite part after merging the large terms}
Since $A_{ab;r}=Z^-_{ab;r}+(2p-m)Z^+_{ab;r}$, define
\begin{equation}
 g_A=g^-,\qquad g_H=g^++(m-2p)g^-,\qquad
 Q_{\core}=\sum(g_HZ^++g_AA_{ab;r}).
\end{equation}
Now impose exact half filling $m=N/2$, $N=p^2/\lambda$, with fixed
$\lambda>0$, and only afterwards consider small $\lambda$.
For fixed nonnegative deficits,
\begin{equation}
 a=p-u,\qquad b=p-v,\qquad r=a+b-S,\qquad q=e^{-4\lambda}.
\end{equation}
The explicit coefficients give
\begin{align}
 \frac{g_A}{C}&=\frac{3-4q+q^2}{2}+O(p^{-1}),\label{eq:ga-limit}\\
 \frac{g_H}{C}&=2p(1-q)^2+K_0
     +8\lambda(q^2-q)(u+v)+8\lambda q^2S+O(p^{-1}),\label{eq:gh-limit}\\
 K_0&=\lambda\left[1-10q+9q^2+16\lambda q-32\lambda q^2\right].
\end{align}
In particular, the order-$m$ terms cancel before the limit is taken.
The remaining order-$p$ coefficient is independent of ranks and
multiplies the conserved operator $H^2-CI$ after centering.

Here is a direct derivation of the finite term. Write
$g_H/C=B_0+B_a\kappa_a/C+B_b\kappa_b/C+B_r\kappa_r/C$ using
$C_-/C=d(n+p)/(mn)$ and $j_c/C=d/m$. For fixed deficits the exact
rational coefficients expand as
\begin{equation}
 (B_0,B_a,B_b,B_r)
 =(2p+\lambda,-2p-5\lambda,-2p-5\lambda,2p+9\lambda)
       +O(p^{-1}).
\end{equation}
For $j=tp-s$, the terminating binomial product gives
\begin{align}
 \frac{\kappa_j}{C}
 &=\prod_{l=0}^{p-1}
       \left(1-\frac{j}{m-l}\right)
       \left(1-\frac{j}{m+p-l}\right),\\
 \log\frac{\kappa_j}{C}
 &=-4\lambda t+\frac{4\lambda s-4\lambda^2t^2}{p}
       +O(p^{-2}).
\end{align}
Taylor's remainder for the logarithm is uniform on any fixed deficit
window. Substituting $(t,s)=(1,u),(1,v),(2,u+v+S)$ gives
Eq.~\eqref{eq:gh-limit}. This is a coefficient expansion; treating
the entire $A$ contribution requires an additional weighted estimate.

\newpage
\section{An $L^2$ result for the $H^2$ contribution}
The $H^2$ part can already be controlled beyond fixed rank windows.
For $a=\alpha p$, $b=\beta p$, $r=\rho p$, the rational coefficients
have the uniform expansions
\begin{align}
 B_0&=2p+\lambda+O(p^{-1}),\\
 B_a&=-2p+\lambda(6-11\alpha)+O(p^{-1}),\\
 B_b&=-2p+\lambda(6-11\beta)+O(p^{-1}),\\
 B_r&=2p+\lambda(5\alpha+5\beta+6\rho-13)+O(p^{-1}).
\end{align}
Uniformity holds for $\alpha,\beta\in[0,1]$, $\rho\in[0,2]$:
after scaling, every rational denominator stays away from zero.
The binomial product and its rank finite differences satisfy
$|\kappa_{j+1}/C-\kappa_j/C|\le K_\lambda/p$ throughout this window.
Comparison with $j=p,2p$ therefore gives
\begin{equation}
 \left|\frac{g_H}{C}-2p(1-q)^2\right|
       \le K_\lambda(1+u+v+S).
 \label{eq:bound}
\end{equation}

Let $C_aI=\EE h_a^2$, $\omega_a=C_a/C$. Different input-rank pairs
are orthogonal in second Gaussian chaos, and output ranks are HS
orthogonal. The previous independent-product law $\pi_{ab}^{(m)}(r)$
and symmetrization bound imply, for any real filter $f$,
\begin{equation}
 \left\|C\sum_{a\le b,r} f_{ab;r}W_2Z^+_{ab;r}\right\|_2^2
 \le4C^4\sum_{a,b}\omega_a\omega_b
                  \sum_r\pi_{ab}^{(m)}(r)|f_{ab;r}|^2.
\end{equation}
The established input-deficit factorial tail and uniform third moment
of $S$ make the square of Eq.~\eqref{eq:bound} uniformly integrable.
First restrict $u,v$, then use the $O(K^{-1})$ squared tail for $S$;
the complementary input tail vanishes by the factorial bound and
the uniform conditional second moment of $S$.

In the physical normalization $C\to J_s^2/\lambda$, define
\begin{equation}
 D_N=\sum_a(p-a)h_a,\qquad
 \mathcal S_N=\sum_{a\le b,r}(a+b-r)W_2Z^+_{ab;r},\qquad
 \mathcal H_N=\sum_{a\le b,r}g_HW_2Z^+_{ab;r}.
\end{equation}
Fixed-window convergence and the preceding tails prove
\begin{align}
 \mathcal H_N={}&C[2p(1-q)^2+K_0]W_2h^2\notag\\
 &+8\lambda C(q^2-q)W_2\{h,D_N\}
       +8\lambda Cq^2\mathcal S_N+o_{L^2}(1).
 \label{eq:l2}
\end{align}
This result is at exact half filling. It does not yet include the
weighted $A$ contribution or the full charge-window average.

\newpage
\section{The remaining spectral problem and verification}
Equation~\eqref{eq:l2} separates a conserved background, an
input-deficit insertion, and an output-deficit insertion. The former
input-filter work identifies the centered linear deficit with the
existing $F$ direction. The new task is the coherent norm and cross
spectrum of $\mathcal S_N$ with $A$ and $J_hF$.

At half filling the exact selection rule supplies a useful starting
point. For independent normalized pure-rank Hermitian $U_a,V_b$,
\begin{equation}
 P_r(V_bU_a)=(-1)^{a+b-r}P_r(U_aV_b),
\end{equation}
and therefore
\begin{equation}
 \EE\|P_r\{U_a,V_b\}\|_m^2
   =4\,\mathbf1_{\{a+b-r\ {
m even}\}}\,
                       \pi_{ab}^{(m)}(r).
\end{equation}
Thus the same-charge symmetric Gram requires parity-weighted
deficit moments. The previous first three ordinary moments do not
alone determine them. The lower-charge and mixed Gram retain their
coherent recoupling data; existing reduced-product formulas
\cite{KV} explicitly retain multiplicity and adjacent-sector phases.

The $A$ coefficient in Eq.~\eqref{eq:ga-limit} is bounded and has a
fixed-deficit limit, but replacing it by that limit under the full
operator sum still needs a weighted tail estimate for $A_{ab;r}$.
The crude bound on $Z^-$ separately loses the cancellation of its
order-$m$ part against $Z^+$. We have not used that crude bound to
claim convergence of the entire memory core. The final raw $Z_2$
also requires restoring the low-chaos terms of the fourth-chaos
descendants before its spectrum can be assembled.

\paragraph{Evidence.}
Five systems, $(N,p,m)=(4,1,2),(5,2,3),(6,2,3),(6,2,4),(6,3,4)$,
give 522 exact rank-projected checks of the five individual
contractions and their core sum on complex, pure-rank inputs.
Separately, Eq.~\eqref{eq:transfer} is checked as a complete linear
map for every external index and output rank: 87 map identities,
5,470,368 complex-entry equalities, maximum cleared integer 60.
These exhaustive transfer checks complement the analytic proof;
the product checks alone are not a full quadratic-tensor audit.

Symbolic algebra independently merges the five coefficients,
derives the finite term and the uniform rational expansions.
Twenty-four finite cases use exact binomial fractions and 80/120
digit evaluations, with actual finite-$p$ deviations retained. For
$\lambda=1/2,p=128,(u,v,S)=(2,1,3)$, the $g_A/C$ deviation is about
$0.02487$; the finite $g_H/C$ deviation is about $-0.02144$.
Precision agreement is not interval certification. All outputs and
the six-page compiled note are saved locally, without Git publication.

\begin{thebibliography}{9}
\bibitem{KV} V. K. B. Kota and M. Vyas,
\emph{Random Matrix Theory for Transition Strength Densities in Finite
Quantum Systems}, Secs.~3, 7, Eqs.~(34)--(35), (113),
\href{https://arxiv.org/abs/1504.00759}{arXiv:1504.00759}.
\end{thebibliography}
\end{document}
